% !TeX program = pdflatex
% ALIUS Bulletin native reconstruction source.
% Compile from the ALIUS-bulletin project root. This file does not include a pre-existing article PDF.
\providecommand{\ALIUSRootPrefix}{}

\ifdefined\ALIUSIssueBuild
\else
\documentclass[a4paper,11pt]{article}
\usepackage[margin=22mm]{geometry}
\usepackage[T1]{fontenc}
\usepackage[utf8]{inputenc}
\usepackage[hidelinks]{hyperref}
\usepackage[protrusion=true,expansion=false]{microtype}
\usepackage{parskip}
\fi
\title{Conceptual, anthropological and cognitive issues surrounding religious experience}
\author{Ann Taves Interviewed by Martin E. Fortier \& Maddalena Canna}
\date{ALIUS Bulletin 2}

\ifdefined\ALIUSIssueBuild
\else
\begin{document}
\fi
\maketitle
\section*{Metadata}
\textbf{Piece type:} interview\par
\textbf{DOI:} 10.34700/3qja-m432\par
\textbf{Keywords:} cognitive science of religion, building-block approach, predictive coding, mystical insight, non-ordinary experience, belief\par
\section*{Abstract}
Cognitive Sciences of Religion scholar Ann Taves is the proponent of a ground-breaking building block approach (BBA) to religious experience. According to Taves, religious experience can be disaggregated into fundamental, constitutive components. Philosopher Martin Fortier and anthropologist Maddalena Canna explore the conceptual, anthropological and cognitive aspects of the foundations of religion, as disaggregated by Taves. In her analyses of the cognitive underpinnings of religion, Taves adopts a Predictive Coding Framework (PCF). The compatibility between PCF and BBA is discussed at the light of the debate opposing inherentist and attributionist theories of religion. Particular attention is given to the role of paradox and non-ordinary experiences (NOE) in the emergence of belief. Finally, other fundamental components of religion are explored, such as the relation between perception (e.g. the vividness of a spiritual vision) and attribution of reality, as well as the interplay between religious processes (e.g. spiritual paths, trajectories) and events (e.g. enlightenment, epiphanies, mystical insights). Empirical examples include case studies from Tantric Buddhism, Kashmir Shivaism, Caribbean Spirit Possession, Amazonian Animism and Ayahuasca ritual.\par
\section*{Extracted Text}
\begingroup
\small
\noindent 109 Conceptual, anthropological and cognitive issues surrounding religious experience An interview with Ann Taves By Martin Fortier \& Maddalena Canna Ann Taves taves@religion.ucsb.edu Department of Religious Studies University of California, Santa Barbara, USA Martin Fortier martin.fortier@ens.fr Institut Jean Nicod ENS/EHESS, Paris, France Maddalena Canna maddalena.canna@ehess.fr Laboratoire d'Anthropologie Sociale EHESS/College de France, Paris, France Citation: Taves, A., Fortier, M. \& Canna, M. (2018). Conceptual, anthropological and cognitive issues surrounding religious experience. An interview with Ann Taves. ALIUS Bulletin, 2, 109-128. A. T. wishes to thank Egil Asprem for his helpful comments on the first draft of this interview. You have proposed to conceive of religion as special paths to a goal (Taves, 2009: 66). If we look at the diversity of religious experience, we notice that, in most cases, the notion of path is associated with a rather conflicting notion of immediate and/or inherent presence of the religious in humans. Various traditions suggest that no path is needed, and that the risk of identifying spiritual practice with a sequence of goal- oriented actions must be avoided by any means. For example, in the Dzogchen tradition, the aspirant can follow a path only under the condition that she is aware that the desired state already pre-exists in herself (Norbu \& Shane 1986). Similar conceptions can be found in some Hindu traditions (e.g., Hughes, 1994). In the Christian world, the debate about the notion of grace, which can be defined as a quality independent from human action, shows how gradualist (path-related) and immediatist conceptions of religious experience are often intertwined and complement one another. In the Buddhist tradition, the debate between subitism and gradualism presents a similar conundrum. In subitism, enlightenment is taken to be attainable all at once, whereas in gradualism, it is said to be achieved only through arduous improvement (Gregory, 1991; Faure, 1991). Similar debates are also present among the different branches of Kashmir Shivaism (Padoux, 2017). Thus, the notion of path seems particularly relevant to account for the gradualist approach to religion, where events play a crucial role. On the other hand, the notion of self-recognition might seem more relevant for the immediatist approach, where A. Taves - Issues surrounding religious experience ALIUS Bulletin n2 (2018) aliusresearch.org/bulletin 110 events play a more peripheral role. How do you conceive of the interplay between gradual and immediate ingredients of religious practice and experience? What is the relation between shifts in self- awareness/self-recognition and special events? Is self-recognition conceivable as a particular kind of special event or do we need slightly different conceptual tools to analyze it? Even more broadly, what does this tension between the gradualist and the immediatist take on religion tell us about the emergence of religion itself? I think the key thing to notice here is the difference between relying on a path schema and claiming that following a path (a gradualist approach) will get you where the tradition says you want to go. The path schema (or something like it, e.g., "a way") seems to play a role in all the examples you give. When I suggested (Taves, 2010, p. 175) that religions, spiritualities, and philosophies are often organized around path schemas, I indicated that they could use the path schema to assess, rank, manipulate, and sometimes transcend things that matter. I didn't elaborate on the transcending idea, but that is what I think is at play in the more immediatist take on paths. It is only by starting on the path with its postulated goal that people ask the question (how do I get there) that the immediatists are answering (by saying they have already arrived). If they don't ask the question, then it is pointless to tell them that they already have the answer. Traditions often encourage people to start out on the path (i.e., do some sort of practice), but warn them that the path will not get them to the goal. Ultimately, the immediatists say you have to transcend the path or recognize that there "is no path" or that it is all about grace and nothing you can do will get you there, etc. To speak to one of your later questions, I think that traditions often use the path schema to set up a paradox, e.g., the path is "no path". Meditating on paradoxes can, I think, trigger shifts in self-awareness, which I would conceptualize as an "internal event". In other words, I think that the tension between gradualist (practice will get you there) and immediatist (only direct insight will get you there) approaches can be used to set up a paradox and that meditation on paradoxes can trigger shifts in self-awareness and deep insights of the sort that immediatists often seek to attain. But people must want to get "there"-even if "there" is the insight that there is no "there"-for either approach to "work". You argue that enhanced vividness of sensory content causes people to take what they experience as more real (Taves, 2009, p. 159-160). Do you consider sensory " Meditating on paradoxes can, I think, trigger shifts in self-awareness, which I would conceptualize as an 'internal event'. " A. Taves - Issues surrounding religious experience ALIUS Bulletin n2 (2018) aliusresearch.org/bulletin 111 vividness and attribution of reality to always be correlated? Tantric visualizations aim at enhancing the vividness of imagination (Kozhevnikov, Louchakova, Josipovic, \& Motes, 2009) and yet their function is precisely to induce a sense of decreased reality (to become aware that any sensation is nothing but the non-real product of the mind) (Beyer, 1973). This seems to be a clear counterexample challenging the putative correlation between vividness and attribution of reality. Another example of dissociation between the sense of reality and the vividness of the sensory content is provided by the comparative study of hallucinations: the sense of reality of psychedelics-induced hallucinations, deliriant-induced hallucinations and Charles Bonnet hallucinations can vary to a great extent, whereas their respective sensory content remains quite stable (Fortier, Forthcoming). Namely, regardless of the vividness of the visual hallucinations, deliriant-induced hallucinations always have a high sense of reality, psychedelics-induced hallucinations a moderate sense of reality and Charles Bonnet hallucinations no sense of reality. According to you, what are the key factors contributing to the attribution of reality to experience in general and to religious experience in particular? I would view vividness and reality as contingent. I think that "paths" assert different conceptions of reality. As animals, I think we have evolved mechanisms (our senses) to determine what is real for us as human animals. Our "natural" or "default" sense of what is real is grounded in the way we have evolved to process input from our bodies and our environment. Various practices and neurological conditions, e.g., Tantric visualizations, psychedelics, or Charles Bonnet syndrome, can all mess with our natural or default sense of what is real. The paths may do this deliberately and then inform us of how we are to appraise what we have experienced. You have recently proposed that the Building-Block Approach (BBA) to religion (Taves, 2009; Taves, 2015) could be fruitfully combined with the Predictive Coding Framework (PCF) (Taves \& Asprem, 2017). Within the PCF, perception and cognition are said to result from a bidirectional tradeoff between top-down predictions and bottom-up prediction errors. Rephrased in the terms of the PCF, the BBA would thus amount to saying that religious experiences and concepts are constructs resulting from the combination of interpretation (top-down processing) and sensory information (bottom-up processing). Let me interrupt the question at this point to clarify our understanding of a few key terms. Building Blocks: When we refer to "building blocks" in the context of the BBA, we are using building blocks as a synonym for the components or parts that interact to A. Taves - Issues surrounding religious experience ALIUS Bulletin n2 (2018) aliusresearch.org/bulletin 112 produce a phenomenon of interest. These components interact to form mechanisms that produce or maintain the phenomenon (Asprem \& Taves, 2016). Mechanisms are nested in stacks and the phenomena of interest can be specified at any given level. Religion: We view "religion" as a complex cultural concept, i.e., as an abstract noun with unstable, overlapping meanings that vary within and across cultures and social formations (Asprem \& Taves, 2016). We refrain from defining religion, so that we can study the way others use it and related terms. Experience: The flow of information in so far as we are aware of it. An experience: an internal event that we have segmented out of the flow of information and thus cognize as "an event." The process of segmenting and cognizing (known as "event cognition") takes place below the threshold of consciousness. We can recount these events and consciously reappraise them after the fact. Narrating and reappraising, however, are new events, which are also cognized subpersonally (Taves \& Asprem, 2017). Attribution theories typically focus on conscious (post-hoc) appraisals of events. One possible worry is that the BBA and the PCF are in fact incompatible. Indeed, the BBA is mainly interested in the cognitive interpretation and appraisal of special events or experiences. It states that a given percept (i.e., a special experience) can be interpreted in various ways depending on one's own expectations and predictions. Now, arguably, this is not what the PCF is mainly concerned with. When the PCF acknowledges the importance of the "interpretation" of sensory data by prior expectations, it means it in a very different way from what the BBA has in mind: interpretation, in the PCF, does not refer to how percepts are cognitively interpreted but how percepts are generated in the first place. In other words, the bidirectional process the PCF is chiefly discussing concerns the production of percepts, whereas the bidirectional process the BBA is focusing on concerns the interpretation of an already-constituted percept. To make this point even clearer, let us consider some examples. Case 1: seeing a coffee cup on the table. Case 2: seeing a face in the foam of a coffee cup. One can be tempted to say that Case 1 and Case 2 both involve some kind of interpretation (some kind of top-down processing). In Case 1, perceptual priors make the brain interpret the sparse sensory data collected by the retina as being a coffee cup. In Case 2, the percept-the foam in the coffee cup-is interpreted as containing a face in virtue of prior experience and gestaltic skills. It must be stressed, however, that two kinds of "interpretation" are here at play: in Case 1, priors contribute to the generation of a percept; in Case 2, priors contribute to the cognitive sense making of a percept which is already there. So, in Case 1, "interpretation" refers to the influence of strictly subpersonal perceptual priors, whereas in Case 2, "interpretation" refers to cognitive priors which may be either personal or subpersonal. A. Taves - Issues surrounding religious experience ALIUS Bulletin n2 (2018) aliusresearch.org/bulletin 113 Importantly, studies looking at the role of priors in the brain mainly focus on cases similar to Case 1 above. For example, when the light-from-above prior makes one see a circle as convex instead of concave, the output of the process is a percept and not the interpretation of a percept (e.g., Sun \& Perona, 1998). By the same token, when, in a binocular rivalry task, the interplay between predictions and prediction errors results in a switch of percept-a house, then a face, then a house, then a face-, the type of top-down processing at work generates distinct percepts and not distinct interpretations of a percept (e.g., Hohwy, Roepstorff, \& Friston, 2008). By contrast, the type of top-down processing the BBA is interested in is not how distinct percepts are generated depending on one's internal models, as much as it is in how a given percept is cognitively interpreted. For example, given that a coffee cup is in front of me, how will my internal models influence my interpretation of the foam of the coffee cup? While the PCF literature dedicated to the study of perceptual priors is extensive (e.g., Knill \& Richards, 1996), to our knowledge, that which is dedicated to the study of top-down processing in the interpretation of sensory stimuli is extremely reduced. Do you agree with the above conceptual distinction between BBA and PCF? If so, do you think the project of unifying BBA and PCF can nonetheless be successfully achieved? I am not as certain as you are that we can make a sharp distinction between "cognitive" and "perceptual" priors and thus between "interpretation" and "perception." This is the issue researchers are discussing under the heading of cognitive penetrability (Zeimbekis \& Raftopoulos, 2015). Some, as you seem to assume, view the formation of percepts as cognitively impenetrable; others disagree. The central question is at what point specific predictions (appraisals derived from cultural schemas) come into play in terms of processing. Language offers an example that likely blurs the distinction between your Case 1 and Case 2. Based on the language(s) we learn as a child and the neuronal pruning that takes place in light of that specific language, there are population specific differences in what we hear (Roepstorff, Niewohner, \& Beck, 2010). Japanese speakers, for example, don't hear the difference between "r" and "l". I would assume that the formation of a percept in this sense (hearing or not hearing the difference) would take place at a lower level of processing than cognizing an event of the sort Asprem and I were discussing. Regardless of how this issue is resolved-and it appears to be a focus of much research-I think that we can integrate the BBA into a predictive processing framework if we think in terms of predictions (aka appraisals) at different levels of mechanisms within a processing hierarchy. In Taves and Asprem (2017), we are considering predictions as appraisals and the error monitoring process as an appraisal process. Schemas, which may be evolved or learned, inform event models, which in turn generate predictions. A. Taves - Issues surrounding religious experience ALIUS Bulletin n2 (2018) aliusresearch.org/bulletin 114 The key thing to recognize is that event cognition is a subpersonal process. We focus on it, because it identifies the components that interact to produce the phenomenon of interest to us as humanists: events that surface to awareness, including internal events, and, thus, the events humans are inclined to turn into narratives of "experiences". But just because the event surfaces to awareness does not mean that we are aware of the components that interact to produce the experience. Your point that cognitive priors, as you call them, may be either personal or subpersonal is an important one. We are assuming that learned (and thus population specific) priors may be internalized to the point where they shape our experience below the threshold of awareness. If some learned priors, such as language, effect changes in the brain such that we cannot perceive things in any other way that would suggest that these learned priors can operate much more deeply than we might suspect. I will be interested in following research that explores the role of population specific differences at those lower levels of processing. I do think that we will be able to unify the BBA and PCF in a more precise fashion at some point, but for now the key point is that processing involves multiple levels of mechanisms and that population specific processes (e.g., cultural schemas) are clearly implicated at subpersonal levels that generate differences in how we appraise events at levels that are not accessible to us. The above distinction between two types of "interpretation" or "top-down processing" has important consequences when it comes to the debate between inherentism and attributionism. Inherentism claims that some experiences are intrinsically religious, whereas attributionism contends that the religious character of experience is derived from some cognitive ascription or attribution. Importantly, the distinction between perceptual priors and cognitive priors leads to the identification of two types of inherentism: - Strong inherentism is the view that the religious character of religious experience does not result from any top-down processing. According to this view, religious experience would be solely determined by bottom-up processing. Prior knowledge would play no role whatsoever in determining the content of experience. Strong inherentism is blatantly inconsistent with PCF. - Weak inherentism acknowledges that experience is partly determined by prior knowledge. According to weak inherentism, experience is the result of a tradeoff between prior expectation and collection of sensory data. This view is perfectly consistent with the PCF. It still qualifies as a version of inherentism because it denies that cognitive priors shape religious experience. It acknowledges that religious experience-be it exteroceptive, sensorimotor, or interoceptive-is determined by priors but these priors are non-cognitive: they determine what an experiential content will be and not how an experiential content will be interpreted. A. Taves - Issues surrounding religious experience ALIUS Bulletin n2 (2018) aliusresearch.org/bulletin 115 - Attributionism assumes that the religious character of experience results from the cognitive interpretation-by System 1 (fast/unconscious thinking) or System 2 (slow/conscious thinking)-of experience. Weak inherentism and attributionism both agree that prior knowledge determines the religious character of religious experience but the type of prior knowledge the former is interested in-i.e., exteroceptive, sensorimotor, and interoceptive priors-differs from the type of prior knowledge the latter is interested in-i.e., cognitive priors. It is worth noting that weak inherentism is compatible with the rejection of perennialism (i.e., the view that religious experience is to a large extent the same in every culture): since the content of experience is determined by non-cognitive priors and since these priors are to a large extent culture-dependent (e.g., Adams, Graf, \& Ernst, 2004), various types of religious experiences can be triggered across cultures (each culture-specific prior triggers a specific type of intrinsically religious experience). Weak inherentism and attributionism are both compatible with the diversity of religious experiences, but for quite distinct reasons. The former implies that the diversity of experiences stems from the plurality of experiential (exteroceptive, proprioceptive and interoceptive) priors, whereas the latter maintains that the diversity of religious experiences stems from the plurality of cognitive priors. In other terms, the enculturation of religious experience can take different forms. Unlike strong inherentism, weak inherentism and attributionism both recognize that culture influences the content of religious experience. But while classical attributionism focuses on the enculturation of religious experience through cognition, weak inherentism emphasizes the possibility for religious experience to be encultured through exteroception, proprioception and interoception. (See summary in Table 1). In your discussion of the theories of religious experience, you convincingly show how strong inherentism is incompatible with the empirical data provided by proponents of the PCF (Taves \& Asprem, 2017). But, it seems, you do not consider the possibility of there being a type of inherentism perfectly compatible with the PCF-namely, weak inherentism. What would be your main objections against weak inherentism? Would you be tempted to endorse a hybrid theory combining weak inherentism and attributionism or to stick to attributionism and reject any form inherentism? I think your terminology captures the subjective distinction between experiences that we consciously appraise and those, such as your experience of a feminine presence (discussed below), that seem to be as real as the cup. But as my comments on the previous question suggest, I don't think that the distinctions work well from a scientific perspective. Part of the problem has to do with terminology generated by various lines of research in sociology (on framing theory (Goffman, 1974; Snow, A. Taves - Issues surrounding religious experience ALIUS Bulletin n2 (2018) aliusresearch.org/bulletin 116 involvement of bottom-up processing? involvement of top-down processing (exteroceptive, proprioceptive and interoceptive priors)? involvement of top-down processing (cognitive priors)? "religious" in virtue of... enculturatuion of religious experience? strong inherentism yes no no sensory (naturalistic version) or extrasensory (supernaturalistic version) data = inherentism no = perennialism weak inherentism yes yes no sensory data and non-cognitive priors = inherentism yes (exteroceptive, proprioceptive, interoceptive) = pluralism attributionism yes (yes) yes sensory data and cognitive priors = attributionism yes (cognitive) = pluralism Table 1: Strong inherentism, weak inherentism and attributionism A. Taves - Issues surrounding religious experience ALIUS Bulletin n2 (2018) aliusresearch.org/bulletin 117 Rochford, Worden, \& Benford, 1986; Johnston \& Noakes, 2005)), social psychology (on attributional theory (Heider, 1958; Malle, 2004; Proudfoot \& Shaver, 1975; Spilka, Shaver, \& Kirkpatrick, 1985)), and emotion and stress research (on appraisal theory (Arnold, 1960; Lazarus \& Folkman, 1984; Scherer, Schorr, \& Johnstone, 2001)). Asprem and I have taken to using the language of appraisal processes as an umbrella term to encompass all these lines of research, which taken together focus on different levels and components involved in processing sensory data. As already indicated, we view the "predictions" generated through predictive processing as subpersonal appraisals and we are open to the possibility that culture in the form of learned content (e.g., language) shapes what we perceive at a very low level of processing. So again, I hesitate to embrace your clearcut distinction between perceptual and cognitive processing. We tend to use cognitive not to refer to conscious (the ordinary language usage) but as a language used to describe all mental processing (as is common among cognitive scientists), in which case the distinction between perceptual and cognitive processing doesn't make sense. Arguments developed in your book (Taves, 2009) as well as by other authors (e.g., Proudfoot, 1985) have forcefully made the case for at least some degree of attributionism. It now seems difficult, and somewhat naive, to defend classical versions of inherentism. Yet, there might arguably be some room for revised (and more sophisticated) versions of inherentism. One of them would be weak inherentism (as defined above). But even that version might still be too strong. Another and more promising version would be a blending of attributionism and weak inherentism. It seems that at least some religious experiences vindicate inherentism. Here is such an example: encounters with the ayahuasca spirit. I (MF) have tried ayahuasca multiple times. While experiencing the effects of the hallucinogenic brew on my brain and body, I was bearing in mind the debates about the nature of religious experience. One thing that struck me was how easy it proved to discriminate between cases in which the experience of the ayahuasca spirit was the result of an attribution and those in which there was no apparent attribution whatsoever. In the first case, I would have some unusual cognitive, bodily or perceptual experience and interpret it with the local Shipibo-Konibo schema: these unusual experiences must be, I inferred, the ayahuasca spirit (oni ibo yoshin) who is " Asprem and I view the 'predictions' generated through predictive processing as subpersonal appraisals and we are open to the possibility that culture in the form of learned content (e.g., language) shapes what we perceive at a very low level of processing. " A. Taves - Issues surrounding religious experience ALIUS Bulletin n2 (2018) aliusresearch.org/bulletin 118 teaching me something, curing my body, showing me images, etc. In those cases, it was obvious that what I was experiencing was only an unusual experience cognitively interpreted as being an ayahuasca spirit. In another context, in another culture, I would certainly have cognitively interpreted the same unusual experience differently (as being God, a demon, an elf, etc.). Here, as far as I can tell, you are simply distinguishing between a conscious appraisal of cues and cues that are appraised below the threshold of consciousness, which is where most such processing takes place. Yes, but importantly, there were other cases in which it seemed to me that no interpretation was involved. To make an analogy: when you are walking in the forest, and you hear a whistle, you may interpret the whistle as the presence of a human person. In such a case, the "perception" of that person requires some cognitive interpretative process. By contrast, when you are having a coffee with someone in a cafe, the experience of that person does not seem to require any interpretation. The experience that you are then having inherently features a human person. Now, this is precisely what I felt on some exceptional occasions: the ayahuasca spirit was there, as obviously as when I am having a coffee with a friend; her presence was completely compelling and undeniable. She was a feminine presence; she was inhabiting an other-world and yet her world was closely interconnected with my everyday world; she was at the same time authoritative, benevolent, mischievous and incredibly witty; somehow, she felt very familiar; and she was exceedingly powerful (I felt she had the power to change all my beliefs with a single snap of the fingers). Perhaps there is nothing very conclusive in such an experience. As you have thoroughly demonstrated, proponents of inherentism have often adduced compelling personal experiences in favor of inherentism, but such personal reports systematically fail to address the possibility that the processes of attribution at play are implicit and therefore subjectively undetectable. One worry with this line of objection, however, is that it would lead one to say that in the meeting with a friend in a cafe, some attribution is also at work. Yes, from a predictive processing perspective, the cognizing of meeting in a cafe event would rely on a "meeting a friend in a cafe" schema, which would rely sub- schemas related to "friend" and "cafe," and on lower level processing (e.g., facial recognition or other cues) to recognize your friend and to navigate your way around the cafe and so on. Now, if attribution is so widespread, then it loses its analytic interest: it misses the key distinction between experiencing spirits and cognitively inferring that spirits are there. This is one of the places where your distinction between experiencing and A. Taves - Issues surrounding religious experience ALIUS Bulletin n2 (2018) aliusresearch.org/bulletin 119 cognitively inferring is quite confusing. I would rephrase the issue as follows: appraisal processes take place at multiple levels of processing, both above and below the threshold of awareness. In order to understand what people experience, we need to specify the level at which appraisals are taking place. When we say we "experience a spirit" rather than "inferring the presence of spirits", we are making a distinction between an experience that surfaces to consciousness more fully formed in the former case than in the latter. If we look at the components that interact below the threshold of consciousness to give rise to "an experience", it is clear (see Figure 2 in Taves \& Asprem, 2017) that learned (event) schemas can operate below the threshold of consciousness to predict what is happening, i.e., that we are interacting with a spirit. If we don't have a schema that can predict what is happening, the experience would surface as a set of cues that we have to consciously interpret, i.e., infer there is a spirit. The inherentist and attributionist models make very straightforward empirical predictions about the nature of the ayahuasca experience. The attributionist insists that interpretations can take place implicitly. As a consequence, people might fail to detect ongoing interpretative processes. The way scholars could detect these implicit interpretations would be by analyzing the phenomenology of these spirit encounters across cultures. If the attributionist is right, then she should find that, in some cultures and contexts, people report the ayahuasca spirit to be feminine while others report it to be masculine; that some report it to be talkative and provider of existential advice while others report it to be secretive; etc. By contrast, the inherentist's prediction would be that some phenomenological features of the ayahuasca spirit are culture-independent: they do not vary across cultures. Ayahuasca has now become a world-wide phenomenon. People are drinking it across the world. In addition to Westerners, many cultures and subcultures are starting to experiment with this hallucinogenic decoction (including, as surprising as it may sound, Iranian Sufis (Rooks, 2014)). Some of these groups of drinkers are not exposed to the narratives circulating about the ayahuasca experience. As such, these groups provide ideal participants to adjudicate the inherentism vs. attributionism debate. If, without being exposed to the emerging ayahuasca world culture and its narratives, these participants' phenomenological descriptions of the ayahuasca spirit perfectly fit with those provided by Amazonian indigenous people and Westerners, then there might be something intrinsically present in those experiences that vindicates some form of inherentism. Preliminary-and yet unpublished-data suggest that the features of the ayahuasca spirit are indeed strikingly consistent across cultures. But much more investigation needs to be done. I think this sounds like a good way to test at what level the percepts are formed and to what extent they are shaped by culture. I think that it is quite possible that ayahuasca and other hallucinogens generate experiences that have some features in common across cultures. I have been following the research with psilocybin where A. Taves - Issues surrounding religious experience ALIUS Bulletin n2 (2018) aliusresearch.org/bulletin 120 there do seem to be features that are commonly elicited, but the researchers are quite clear that set and setting play an important role in how the experience itself unfolds in real time (Griffiths et al., 2006, 2011). (In fact, Timothy Leary coined the set and setting distinction in the context of the Harvard Psilocybin Project in the early sixties (Hartogsohn, 2017; Zinberg, 1986).) The importance of set and setting is particularly apparent if we recognize that researchers are using psilocybin to model both "mystical" and "psychotic" experiences (see, e.g., Gonzalez-Maeso \& Sealfon, 2009; Vollenweider \& Kometer, 2010)! In the former case, the researchers go to great lengths to produce a positive experience by controlling set and setting. Insofar as they limit themselves to administering the mysticism scales (e.g., the MEQ 30), they only ask participants to report on those aspects of the experience that fit with it the scale's definition of "mystical." If proven right, the form of inherentism sketched above would not vindicate perennialism. Indeed, it would imply that some brain patterns can lead one to perceive a specific spirit presence, but that another-induced, for example, by another hallucinogenic chemical compound-would trigger a different kind of religious experience (e.g., the presence of a spirit exhibiting other features). This revised version of inherentism would reject the classical neurotheological quest for a single "God-spot" (e.g., D'Aquili \& Newberg, 1999). According the revamped version of inherentism, there would be dozens of spirit experiences mapping with dozens of distinct neural correlates-the putative single "God-spot" would thereby be replaced by various "spirits-patterns" or "gods-patterns". Within this hypothetical model, the diversity of religious experiences would be duly acknowledged, but it would remain that the phenomenological features-and the neurobiological signature-of each of these religious experiences would be invariant across cultures-making them intrinsically religious and not religious in virtue of some attribution. Although I think that it is possible that experiences of presence are triggered by some hallucinogens, I do not think you can refer to them as "religious experiences", since experiences of presence can be appraised in many different ways (see Barnby \& Bell, 2017). The determination of who the presence is and the nature of the interaction is, I would think, very much determined by set and setting (as well as lower level appraisal processes that may tend to skew the experience in a positive or negative direction (Underwood, Kumari, \& Peters, 2016a, 2016b)). In your own case, you report a "feminine presence" with certain attributes, which you then characterize as the ayahuasca spirit; this latter characterization is surely a conscious appraisal. The experience-phenomenologically speaking-is of a particular sort of presence that fits with your expectations regarding the "ayahuasca spirit". If naive subjects were given ayahuasca without knowing what drug they were taking in a neutral (or medical) environment, I suspect that their experience of a presence, if they had one, might be quite different and definitely would not be recognized as the A. Taves - Issues surrounding religious experience ALIUS Bulletin n2 (2018) aliusresearch.org/bulletin 121 "ayahuasca spirit". Yes, but if compelling evidence was to support the existence of invariant phenomenological features of the "ayahuasca spirit" across cultures, could it somewhat challenge the attributionist thesis? I am making a fairly subtle distinction here, which you may or may not find meaningful, between what is experienced phenomenologically and how it is consciously appraised. I think it is possible that ayahuasca sometimes triggers cues that are appraised below the threshold of consciousness as a sensed presence. As indicated above, I do not think that naive users would recognize the presence as the "ayahuasca spirit." Their conscious appraisals would reflect their own prior experience and the setting in which they received the drug. You have suggested that events triggering disruption of the ordinary sense of reality foster the search for new meanings and predispose subjects to adopt new religious beliefs (Taves 2009, p. 100). Paradox is by definition a cognitive contradiction which, to be managed, requires one to reframe the context of its appraisal. Thus, paradox works as a trigger of new cognitive solutions. Data collected on my own (MC) fieldsite tend to confirm the link between paradoxical mental activity and Altered States of Consciousness (ASCs) (deemed) spiritual. Mental paradoxical practices among Miskitos lead to dissociative hallucinatory seizures attributed to spirit attacks (Canna, 2017). Most Miskitos experiencing spirit attacks have been exposed to contradictory injunctions. They are told "not to think" about the spirit and yet at the same time are incited to pay attention to any possible manifestation of Him (kaiki bas, lukpara, literally "be careful not to care"). This insoluble contradiction triggers a paradoxical state, dissociative symptoms (trance, somatizations, conversions), and hallucinations. In most cases, the mental image of the spirit is appraised as impossible to "expel from one's head" (in the words of my interlocutors) and this loss of control is re-appraised as a proof of the autonomy of the image, and as a consequence, of the spirit's reality. These data are consistent with William Barnard's self-account of his spiritual experience (Barnard, 1997; Taves, 2009, pp. 102-119). As you underlined, Barnard's experience was triggered by paradoxical mental activity. First, Barnard become obsessed with the idea of what would happen to him after death. He tried to imagine himself as a not existing being, and thus experienced a paradox described as follows: "I kept trying, without success, to envision a simple blank nothingness (the visualization exercise generated a paradox-that is, asking self to imagine self not being able to imagine)" (in Taves, 2009, p. 109). This paradox fostered the emergence of an ASC: "Suddenly, without warning, something shifted inside. I felt lifted outside myself, as if I had been expanded beyond my previous sense of self" (in Taves, 2009, p. 110). A. Taves - Issues surrounding religious experience ALIUS Bulletin n2 (2018) aliusresearch.org/bulletin 122 To what extent do you think that the notion of paradox can explain the emergence of different states of consciousness (deemed) spiritual? Paradox is also pervasively present in many religious traditions (e.g., Christian notion of Trinity). To what extent do you think paradox is a typical feature of religious experience? I think your findings and insights on this are very intriguing. Rather than say that paradox is a typical feature of religious experience, I would prefer to say that there are a number of traditions that use paradox to generate the insights that they claim reveal Truth or Reality (deliberately capitalized). I was thinking about this when I answered the paths question above (see first question). On a more methodological note, you have championed a highly directive method for conducting post-hoc elicitation interviews (Taves, 2009; Taves \& Asprem, 2017). For example, you propose to interrupt the narrative flow by asking the narrator w- questions such as "what" and "why", in order to distinguish the event model she has in mind (i.e., her sense of "what happened") from the explanation she provides (i.e., her explanation as to "why it happened"). Directive interviewing methods have both upsides and downsides. A possible downside is that by interrupting the interviewee's narration, directive methods lose some valuable information: it prevents the interviewee to re-enact her experience and access pieces of experience that she would not be able to recollect otherwise (Petitmengin \& Lachaux, 2013; Petitmengin, Remillieux, Cahour, \& Carter-Thomas, 2013). Very directive and less directive methods tend to produce substantially different data. What is your take on the advantages and disadvantages of the distinct kinds of interviewing method? Asprem and I were writing as historians not ethnographers, so we weren't actually offering a method for conducting "post-hoc elicitation interviews", but rather offering a method for analyzing narratives left to us by dead people, where interrupting the flow the narrative is not an issue! I think your observation is spot on when it comes to interviews of living people and, as far as I am aware, Claire Petitmengin has developed the most sophisticated method for teasing out the nuances of experiences from living subjects. You have suggested that some cognitive impairments make a subject more prone to religious experience (Taves \& Asprem, 2017). Studies by Tanya Luhrmann et al. (2015) as well as Rebecca Seligman and Lawrence Kirmayer (2008) suggest that phenomenologically similar experiences can be experienced as pathological or not depending on the social context. If we agree with these authors, a crucial anthropological question is to establish to what extent the opposition between pathology and health is arbitrary. Namely, boundaries between pathological and altered states of consciousness (deemed) religious seem particularly porous. A. Taves - Issues surrounding religious experience ALIUS Bulletin n2 (2018) aliusresearch.org/bulletin 123 Suffering does not seem to be a distinctive feature of pathological experiences as opposed to religious ones. Indeed, in several shamanic traditions, the shaman's relationship to spirits is experienced as a form of painful oppression or afflicting persecution (Stepanoff, 2015; Canna, 2017). The possibility of controlling spiritual manifestations does not seem to be a critical feature of pathology either (e.g., Halloy, 2015). To what extent does the adoption of socially shared frameworks make anomalous experiences acceptable and non-pathological? How do you conceive of the boundaries between pathology, normality and election? This is a great question and one that requires further research. Appraisal processes clearly make a difference, but it is not clear how much of a difference and for whom. Psychosis researchers are working to answer this precise question (see the work of Underwood, Kumari, \& Peters, 2016a, 2016b). While some scholars-classical (Tylor, 1871) or contemporary (Bulkeley, 2016)- argue that religion originates in anomalous experience such as dreaming, proponents of the cognitive science of religion contend that representations produced by our intuitive ontologies-rather than experience-are sufficient to explain the emergence of religion (e.g., Boyer, 2001). By insisting that there are no intrinsically religious experiences and that experience always has to be processed and interpreted by cognitive schemas to be deemed religious, you seem to reject the classical Tylorian view: experience alone is not sufficient for religion to emerge. Now, importantly, there are two ways of rejecting the hypothesis that the origin of religion lies in experience: (1) Weak rejection: experience is not sufficient but is necessary for religion to emerge. (2) Strong rejection: experience is not sufficient nor necessary for religion to emerge. What is your view on the origin of religion? Do you think that the strong rejection of the Tylorian view is a tenable position? Alternatively, do you lean towards the view that experience is not sufficient but indeed necessary? The answer to this depends on what we mean by "religion" and "experience". If we define "experience" as an internal event that segmented out of the flow of " I think we can assume that 'religion' and other innovations rely on some sort of new idea or insight that provides some sort of competitive advantage. In that sense, all innovations rely on some sort of experience. " A. Taves - Issues surrounding religious experience ALIUS Bulletin n2 (2018) aliusresearch.org/bulletin 124 information of which we are aware, then any idea or insight is an "experience". If we assume that, whatever we mean by "religion" and however many times and places it emerges, it is-as an emergent phenomenon-an innovation. Considered as such, I think we can assume that "religion" and other innovations rely (at a minimum) on some sort of new idea or insight that provides some sort of competitive advantage. In that sense, all innovations rely on some sort of experience (aka insight or idea). A. Taves - Issues surrounding religious experience ALIUS Bulletin n2 (2018) aliusresearch.org/bulletin 125 References Adams, W., Graf, E., \& Ernst, M. (2004). Experience can change the "light-from-above" prior. Nature Neuroscience, 7(10), 1057-1058. https://doi.org/10.1038/nn1312 Arnold, M. (1960). Emotion \& Personality. Columbia University Press. Asprem, E., \& Taves, A. (2016). Building Blocks of Human Experience Website. Retrieved from bbhe.ucsb.edu Barnard, G.W. (1997). Exploring unseen worlds: William James and the Philosophy of Mysticism. Albany: State University of New York Press. Barnby, J. M., \& Bell, V. (2017). The Sensed Presence Questionnaire (SenPQ): initial psychometric validation of a measure of the "Sensed Presence" experience. PeerJ, 5, e3149. https://doi.org/10.7717/peerj.3149 Beyer, S. (1973). Magic and ritual in Tibet: The cult of Tara. Delhi: Motilal Barnasidass. Boyer, P. (2001). Religion Explained: The Evolutionary Origins of Religious Thought. New York: Basic Books. Bulkeley, K. (2016). Big dreams: The science of dreaming and the origins of religion. Oxford/New York: Oxford University Press. Canna, M (2017). Modeliser les etats modifies de conscience. Vers une anthropologie interactionnelle de la conscience. Intellectica, 67, 268-299. D'Aquili, E., \& Newberg, A. (1999). The Mystical Mind: Probing the Biology of Religious Experience. Minneapolis: Fortress Press. Faure, Bernard (1991). The Rhetoric of Immediacy: A Cultural Critique of Chan/Zen Buddhism. Princeton NJ/Oxford: Princeton University Press. Hughes, John. (1994). Self Realization in Kashmir Shivaism: The oral Teachings of Swami Lakshmanjoo. Albany: State University of New York Press. Fortier, M. (Forthcoming). Sense of reality, metacognition and culture in schizophrenic and drug-induced hallucinations: An interdisciplinary approach. In J. Proust \& M. Fortier (Eds.), Metacognitive Diversity: An Interdisciplinary Approach. Oxford/New York: Oxford University Press. Gregory, Peter N. (1991). Sudden and Gradual: Approaches to Enlightenment in Chinese Thought. Motilal Banarsidass Publications. Goffman, E. (1974). Frame Analysis: An Essay on the Organization of Experience. Lebanon, NH: Northeastern University Press. A. Taves - Issues surrounding religious experience ALIUS Bulletin n2 (2018) aliusresearch.org/bulletin 126 Gonzalez-Maeso, J., \& Sealfon, S. C. (2009). Psychedelics and schizophrenia. Trends in Neurosciences, 32(4), 225-232. Griffiths, R. R., Johnson, M. W., Richards, W. A., Richards, B. D., McCann, U., \& Jesse, R. (2011). Psilocybin occasioned mystical-type experiences: immediate and persisting dose-related effects. Psychopharmacology (Berl), 218(4), 649-65. https://doi.org/10.1007/s00213-011-2358-5 Griffiths, R. R., Richards, W. A., McCann, U., \& Jesse, R. (2006). Psilocybin can occasion mystical-type experiences having substantial and sustained personal meaning and spiritual significance. Psychopharmacology (Berl), 187(3), 268-83; discussion 284-92. https://doi.org/10.1007/s00213-006-0457-5 Halloy, A. (2015). Divinites incarnees: L'apprentissage de la possession dans un culte afro-bresilien. Paris: PETRA. Hartogsohn, I. (2017). Constructing drug effects: A history of set and setting. Drug Science, Policy and Law, 3, 1-17. https://doi.org/10.1177/2050324516683325 Heider, F. (1958). The Psychology of Interpersonal Relations. Hillsdale, NJ: Lawrence Erlbaum Associates. Hohwy, J., Roepstorff, A., \& Friston, K. (2008). Predictive coding explains binocular rivalry: an epistemological review. Cognition, 108(3), 687-701. https://doi.org/10.1016/j.cognition.2008.05.010 Johnston, H., \& Noakes, J. A. (2005). Frames of protest: social movements and the framing perspective. Rowman \& Littlefield Publishers. Knill, D., \& Richards, W. (Eds.). (1996). Perception as Bayesian inference. Cambridge: Cambridge University Press. Kozhevnikov, M., Louchakova, O., Josipovic, Z., \& Motes, M. A. (2009). The enhancement of visuospatial processing efficiency through buddhist deity meditation. Psychological Science, 20(5), 645-653. Lazarus, R. S., \& Folkman, S. (1984). Stress, Appraisal, and Coping (1 edition). New York: Springer Publishing Company. Luhrmann, T., Padmavati, R., Tharoor, H., \& Osei, A. (2015). Differences in voice-hearing experiences of people with psychosis in the USA, India and Ghana: interview-based study. The British Journal of Psychiatry, 206(1), 41-44. Malle, B. F. (2004). How the Mind Explains Behavior: Folk Explanations, Meaning, and Social Interaction. Cambridge, Mass: A Bradford Book. Norbu, Namkhai Chogyal \& Shane, John. (1986). The Crystal and the Way of Light: Sutra, A. Taves - Issues surrounding religious experience ALIUS Bulletin n2 (2018) aliusresearch.org/bulletin 127 Tantra, and Dzogchen. New York/London: Routledge \& Kegan Paul. Padoux, A. (2017). The Hindu Tantric World: An Overview. Chicago: Chicago University Press. Petitmengin, C., \& Lachaux, J.-P. (2013). Microcognitive science: Bridging experiential and neuronal microdynamics. Frontiers in Human Neuroscience, 7. https://doi.org/10.3389/fnhum.2013.00617 Petitmengin, C., Remillieux, A., Cahour, B., \& Carter-Thomas, S. (2013). A gap in Nisbett and Wilson's findings? A first-person access to our cognitive processes. Consciousness and Cognition, 22(2), 654-669. https://doi.org/10.1016/j.concog.2013.02.004 Proudfoot, W., \& Shaver, P. (1975). Attribution Theory and the Psychology of Religion. Journal for the Scientific Study of Religion, 14(4), 317. https://doi.org/10.2307/1384404 Roepstorff, A., Niewohner, J., \& Beck, S. (2010). Enculturing brains through patterned practices. Neural Networks, 23(8-9), 1051-1059. https://doi.org/10.1016/j.neunet.2010.08.002 Rooks, B. (2014). Ayahuasca and the Godhead: An Interview with Wahid Azal of the Fatimiya Sufi Order. Reality Sandwich. Retrieved from http://realitysandwich.com/219826/ayahuasca-and-the-godhead-an-interview- with-wahid-azal-of-the-the-fatimiya-sufi-order/ Scherer, K. R., Schorr, A., \& Johnstone, T. (Eds.). (2001). Appraisal Processes in Emotion: Theory, Methods, Research (1 edition). Oxford, New York: Oxford University Press. Seligman, R., Kirmayer, L., (2008). Dissociative Experience and Cultural Neuroscience: Narrative, Metaphor and Mechanism. Cultural Medical Psychiatry, 32, 31-64. Snow, D. A., Rochford, E. B., Worden, S. K., \& Benford, R. D. (1986). Frame Alignment Processes, Micromobilization, and Movement Participation. American Sociological Review, 51(4), 464. https://doi.org/10.2307/2095581 Spilka, B. (B), Shaver, P. (P), \& Kirkpatrick, L. A. (LA). (1985). A general attribution theory for the psychology of religion. Journal for the Scientific Study of Religion, 24(1), 1-20. Stepanoff, Charles. (2015). Trans-singularities: The Cognitive Foundations of Shamanism in Northern Asia. Social Anthropology, 23 (2), 169-185. Sun, J., \& Perona, P. (1998). Where is the sun? Nature Neuroscience, 1(3), 183-184. https://doi.org/10.1038/630 Taves, A. (2009). Religious experience reconsidered: A building-block approach to the study of religion and other special things. Princeton NJ/Oxford: Princeton University Press. A. Taves - Issues surrounding religious experience ALIUS Bulletin n2 (2018) aliusresearch.org/bulletin 128 Taves, A. (2010). No Field Is an Island: Fostering Collaboration between the Academic Study of Religion and the Sciences. Method and Theory in the Study of Religion, 22(2- 3), 177-188. Taves, A. (2015). Reverse Engineering Complex Cultural Concepts: Identifying Building Blocks of "Religion." Journal of Cognition and Culture, 15(1-2), 191-216. https://doi.org/10.1163/15685373-12342146 Taves, A., \& Asprem, E. (2017). Experience as event: Event cognition and the study of (religious) experiences. Religion, Brain \& Behavior, 7(1), 43-62. https://doi.org/10.1080/2153599X.2016.1150327 Tylor, E. (1871). Primitive culture: Researches into the development of mythology, philosophy, religion, language, art and custom. Londres: John Murray. Underwood, R., Kumari, V., \& Peters, E. (2016). Appraisals of psychotic experiences: an experimental investigation of symptomatic, remitted and non-need-for-care individuals. Psychol Med, 46(6), 1249-63. https://doi.org/10.1017/S0033291715002780 Underwood, R., Kumari, V., \& Peters, E. (2016). Cognitive and neural models of threat appraisal in psychosis: A theoretical integration. Psychiatry Res, 239, 131-8. https://doi.org/10.1016/j.psychres.2016.03.016 Vollenweider, F. X., \& Kometer, M. (2010). The neurobiology of psychedelic drugs: implications for the treatment of mood disorders. Nature Reviews Neuroscience, 11(9), 642-651. Zeimbekis, J., \& Raftopoulos, A. (Eds.). (2015). The Cognitive Penetrability of Perception: New Philosophical Perspectives (1 edition). New York, NY: Oxford University Press. Zinberg, N. (1986). Drug, Set, and Setting: The Basis for Controlled Intoxicant Use (1 edition). New Haven: Yale University Press.\par
\endgroup

\ifdefined\ALIUSIssueBuild
\clearpage
\expandafter\endinput
\fi
\end{document}
